% page 318 du fac-simile (image p-317 du scan)
% bloc 152.4 x 242.0 mm ; marge gauche 8.0 mm
\pgc{-12.700mm}{0.000mm}{
\l{\cel{3}310}
\l{}
\l{}
\l{\sou{kosmografio}. Deskripto di la sistemo astronomiala di nia}
\l{\cel{3}planeto e di universo. DEFIRS.}
\l{}
\l{\sou{kosmologio}. La cienco pri la +leyi (legi) qui regnas}
\l{\cel{3}universo. - DEFIRS.}
\l{}
\l{\sou{kosmopolita.}\cel{2}I. Qua egardas su kom civitano di mondo}
\l{\cel{3}(kontraste kun civitano di ca o ta lando). -}
\l{\cel{3}II. Qua vivas lore en ca lando, lore en ita, lore en}
\l{\cel{3}altra, senfine. - DEFIRS.}
\l{}
\l{\sou{kosto}. (\sou{anat.)}\cel{3}Singla ek la osti plata ed arkatra qui,}
\l{\cel{3}departante de la spino, e venante rajuntar avane la}
\l{\cel{3}sternumo, formacas la kavajo ostoza \textquotedbl{}pektoro\textquotedbl{}. - EFIS.}
\l{}
\l{\sou{kostumo}. Maniero vestizar su, propra di populo, di epoko,}
\l{\cel{3}di klaso, di cirkonstanco. - DEFIR.}
\l{}
\l{\sou{koto}. I. Indiko pri la altitudo di la punti diversa di}
\l{\cel{3}tereno, di aquo-fluo, e c., relate nivelo egardata kom}
\l{\cel{3}bazo. - II. Indiko, per cifri, per literi, di singla}
\l{\cel{3}serio di peci ek dokumentaro, ek arkiveyo. - F.}
\l{}
\l{\sou{koteno.}\cel{2}(\sou{patol.})\cel{2}Falsa membrano qua formacesas interne}
\l{\cel{3}di la guturo, en ula kazi di angino. - I.}
\l{}
\l{\sou{kotiledono.}\cel{2}I. (\sou{anat.)}\cel{2}Lobo karnoza di la placenta. -}
\l{\cel{3}II. (\sou{bot.}) Lobo, folio embriona di la semino di veje-}
\l{\cel{3}tanto fanerogama. - DEFIS.}
\l{}
\l{\sou{kotiliono}.\cel{2}Danso kun figuri qua finas la balo (dansata}
\l{\cel{3}olim da 4 od 8 personi). - DEFIS.}
\l{}
\l{\sou{kotiloido}. (\sou{anat.}) Di qua la formo esas olta di ta kavajo}
\l{\cel{3}di osto qua recevas altra osto. -}
\l{}
\l{\sou{kotleto}. (\sou{koquarto)}.\cel{2}Kosto de mutono, de bovyuno, de por-}
\l{\cel{3}ko, ec., destinita a manjeso da homo. - DEFIR.}
\l{}
\l{\sou{kotono}. Materio texebla, blanka, fina, qua kovras la semino}
\l{\cel{3}di arboreto qua formacas genero ek la familio \textquotedbl{}malvacei\textquotedbl{},}
\l{\cel{3}qua origine venis de India e quan onu kultivas nun anke}
\l{\cel{3}en Egiptia ed en la regiono sud-esta di Usa. - EFIS.}
\l{}
\l{\sou{koturno}. Pedovesto quan +weris, che la populi antiqua, la}
\l{\cel{3}viri e la mulieri di rango super-meza ; sandalo kun}
\l{\cel{3}suolo, fixigita an la pedo per ledro-bendi ligita super}
\l{\cel{3}la maleolo e lasante la pedo-fingri nekovrata. - DFIS.}
\l{}
\l{\sou{kovar}. (trans.)\cel{2}(\sou{aludante la femino di ucelo)}.\cel{2}Sternar}
\l{\cel{3}su e permanar tale dum ula quanto de tempo, sur sua ovi,}
\l{\cel{3}por konservar ad oli konstante la kaloro qua igos la}
\l{\cel{3}uceleto ekirar kande parformacita. - FI.}
}
